\section{Prompt Outlines}
\label{appendix:prompt-outlines}

This appendix summarizes the prompt outlines used in LMVQA. The outlines abstract the implemented prompts by showing their roles, inputs, grounding rules, and expected output formats.

\newcommand{\LmvqaPromptBox}[2]{%
\vspace{0.55em}
\noindent\fbox{%
\begin{minipage}{0.96\linewidth}
\small
\textbf{#1}\\[0.25em]
#2
\end{minipage}}
\vspace{0.35em}
}

\newcommand{\LmvqaPItem}[2]{\textbf{#1} #2\\}

\subsection{Main LMVQA Prompt Outlines}
\label{appendix:main-prompt-outlines}

\LmvqaPromptBox{(a) Prompt Outline for Non-Diagram Frame Captioning}{
\LmvqaPItem{(I)}{Role: an assistant analyzing a video frame from an educational lecture or tutorial.}
\LmvqaPItem{(II)}{Frame metadata: \texttt{\{timestamp\}}.}
\LmvqaPItem{(III)}{Input: a sampled video frame image.}
\LmvqaPItem{(IV)}{Task: describe visible slide, screen, chart, or non-diagram content; extract and summarize visible text when useful.}
\LmvqaPItem{(V)}{Grounding rules: ignore irrelevant UI elements, avoid redundancy, and describe only visible content.}
\LmvqaPItem{(VI)}{Output: a fluent structured English frame description.}
}

\LmvqaPromptBox{(b) Prompt Outline for Diagram-Like Frame Detection}{
\LmvqaPItem{(I)}{Role: a fast triage classifier.}
\LmvqaPItem{(II)}{Input: a single video frame image.}
\LmvqaPItem{(III)}{Task: decide whether the frame is diagram-like. Diagram-like visuals include diagrams, flowcharts, schematics, architecture diagrams, charts, tables, UI wireframes, whiteboards with boxes or arrows, and equations.}
\LmvqaPItem{(IV)}{Output: \texttt{LABEL: <DIAGRAM|NOT\_DIAGRAM>}; \texttt{TYPES: <types>}; \texttt{CONF: <0.00-1.00>}.}
}

\LmvqaPromptBox{(c) Prompt Outline for Seven-Way Diagram Classification}{
\LmvqaPItem{(I)}{Role: a fast diagram-type classifier.}
\LmvqaPItem{(II)}{Input condition: the image is already known to be diagram-like.}
\LmvqaPItem{(III)}{Categories: \texttt{UML\_CLASS}, \texttt{UML\_SEQUENCE}, \texttt{UML\_STATE\_ACTIVITY}, \texttt{UML\_USE\_CASE}, \texttt{NETWORK\_TOPOLOGY}, \texttt{ARCH\_WORKFLOW}, or \texttt{OTHER}.}
\LmvqaPItem{(IV)}{Output: \texttt{CATEGORY: <category>}; \texttt{CONF: <0.00-1.00>}.}
}

\LmvqaPromptBox{(d) Seven Diagram Categories Used for Category-Specific Extraction}{
\LmvqaPItem{\texttt{UML\_CLASS}:}{diagram title; scope or packages; classes and interfaces; relationships; key takeaways.}
\LmvqaPItem{\texttt{UML\_SEQUENCE}:}{diagram title; participants or lifelines; temporally ordered messages; activations; combined fragments; key scenario summary.}
\LmvqaPItem{\texttt{UML\_STATE\_ACTIVITY}:}{diagram title; subtype; nodes; transitions or flows; composite states, swimlanes, forks, joins, and behavior summary.}
\LmvqaPItem{\texttt{UML\_USE\_CASE}:}{diagram title; system boundary; actors; use cases; associations; include, extend, and generalization relationships.}
\LmvqaPItem{\texttt{NETWORK\_TOPOLOGY}:}{title or legend; nodes and endpoints; links; traffic or flows; key topology takeaway.}
\LmvqaPItem{\texttt{ARCH\_WORKFLOW}:}{title; components or modules; interfaces or artifacts; connections or flows; workflow steps.}
\LmvqaPItem{\texttt{OTHER}:}{title; diagram type guess; main elements; relationships or structure; key takeaway.}
}

\LmvqaPromptBox{(e) Prompt Outline for Diagram Frame Extraction}{
\LmvqaPItem{(I)}{Role: a diagram-frame describer specialized according to the selected diagram category.}
\LmvqaPItem{(II)}{Frame metadata: \texttt{\{timestamp\}}.}
\LmvqaPItem{(III)}{Input: the diagram-like video frame image.}
\LmvqaPItem{(IV)}{Generic schema: title, purpose, global layout, nodes, edges, flow sequence, text and labels, data or variables, context anchors, and key takeaway.}
\LmvqaPItem{(V)}{Category-specific schema: refined according to the seven-way classifier output.}
\LmvqaPItem{(VI)}{Grounding rules: describe only visible content, quote labels when possible, mark unclear text as \texttt{[illegible]}, and do not invent facts.}
\LmvqaPItem{(VII)}{Output: exhaustive structured plain English text, not JSON.}
}

\LmvqaPromptBox{(f) Prompt Outline for RAG-Based Video Question Answering}{
\LmvqaPItem{(I)}{Role: an assistant answering questions about lecture videos using timestamped multimodal information.}
\LmvqaPItem{(II)}{Video metadata: \texttt{\{video\_duration\}}.}
\LmvqaPItem{(III)}{Retrieved context: \texttt{\{retrieved\_context\}}, formatted as \texttt{[source] <timestamp>: <text>}.}
\LmvqaPItem{(IV)}{User question: \texttt{\{user\_question\}}.}
\LmvqaPItem{(V)}{Grounding rules: answer only from retrieved segments; do not introduce unsupported information; use the insufficient-information answer when needed.}
\LmvqaPItem{(VI)}{Output: \texttt{Answer:} plus final answer, and \texttt{Used\_context:} with 1--10 relevant timestamped segments.}
}

\begin{center}
\small\textbf{Fig. 1: Outlines of the main prompts used in LMVQA.}
\end{center}
